\label{chap:83-02}

\section{Geometric realizations of TRIT: elliptic, parabolic and hyperbolic regimes and double temporal orientation}

TRIT preserves within a single state the published residue, the carry
cocycle, orientation, and compositional memory. Its \(3\) algebraic
regimes are not editorial layers: they are correlative realizations of the
quadratic operator \(J_{\tau}^{2}=-\tau I\), linked to the \(2\)
orientations of time and to the threshold state.

\begin{HMTScientificOwnerSubordinated}
\def\HMTTRITFormal{1}
\def\HMTTRITDevelopments{1}
\input{sections/hmt/02_trit_geometrias}
\let\HMTTRITFormal\undefined
\let\HMTTRITDevelopments\undefined
\end{HMTScientificOwnerSubordinated}

% RESULTADO_RECUPERADO. Owner integro: sections/hmt/09b_algebras_cuadraticas.tex
% Destino causal: capitulo 2; realizaciones cuadraticas y generadores TRIT.
\begin{HMTScientificOwnerSubordinated}
\input{sections/hmt/09b_algebras_cuadraticas}
\end{HMTScientificOwnerSubordinated}

% RESULTADO_RECUPERADO. Owner integro: sections/hmt/03k_genealogia_operaciones_rev2.tex
% Destino causal: capitulo 2; operaciones levantadas y memoria del TPK.
\begin{HMTScientificOwnerSubordinated}
\input{sections/hmt/03k_genealogia_operaciones_rev2}
\end{HMTScientificOwnerSubordinated}

\section[APP, TRIT and TPK: support, orientation and transport]
{APP, TRIT and TPK: support, orientation and transport}

\subsection{APP as fiber arithmetic}

APP is not an illustrative table: it is the arithmetic support on which the
difference between aggregation and composition becomes visible. In
\(
X_9=(\Znine)^2
\),
summation accumulates displacements and alternative histories; multiplication composes
stages and, when a factor is not invertible, folds distinct inputs onto
the same residue. The residue--quotient lift unfolds those fibers and
preserves the genealogy contracted by the modular publication.

Let \(q(a,b)=a+b\pmod 9\). The published datum \(q(a,b)\) does not determine the pair
\((a,b)\). The fiber
\[
 q^{-1}(r)=\{(a,b):a+b\equiv r\pmod9\}
\]
is memory of the alternatives. For multiplication, the contraction is even
more pronounced: the nonunit classes \((0,3,6)\) produce fibers of distinct
sizes. APP therefore publishes a geometry of operations before
any physical interpretation.

The active central block is
\[
B_c=\begin{pmatrix}7&2\\2&7\end{pmatrix}
   =7I+2S=9P_++5P_-,
\qquad
P_\pm=\frac{I\pm S}{2}.
\]
Its local spectral jump is \(9-5=4\); the coupling coefficient is \(2\).
The value \(54\) belongs to a later unfolding and, in the electronic fiber,
to \(D_1(\Gamma_e)\). Equating them because they occur in a numerical chain would erase
their types.

\subsection{Local Lorentz geometry and Minkowski plane \texorpdfstring{\(1+1\)}{1+1}}

Normalization
\[
M_c=\frac{B_c}{\sqrt{45}}
=\exp\!\left(\frac12\log\frac95\,S\right)
\]
realizes an element of \(SO^+(1,1)\). APP thus contains, prior to continuous
physics, a decomposition into \(P_+,P_-\) modes, a gap, a coupling,
and a velocity. The result is not yet Minkowski space-time: it is
the algebraic structure that will make a Lorentzian representation possible.

\subsection{TRIT as orientation and local memory}

TRIT introduces the sign that the modular publication omits. In a triadic sum.
\[
 a+b=r+3c,
\]
the residue \(r\) publishes the result, and the carry \(c\) preserves how it was obtained. The states \((+1,0,-1)\) should not be interpreted solely as \(3\) colors. In the temporal parametrization they act, respectively, as return with memory of the past, present threshold, and bidirectional opening toward the future; in the algebraic realization they distinguish counterflow, fixation, and flow; in the geometric realization they label the hyperbolic, parabolic, and elliptic branches of a quadratic family. Compatibility among these readings requires maps, not a verbal identification.

\begin{definition}[Enhanced TRIT state]
A TRIT state is a pair \((r,c)\) endowed with orientation. The first
component is publishable; the second belongs to the memory fiber. Change of
reader may forget \(c\), but TPK transport cannot do so.
\end{definition}

This separation makes it possible to formalize present, past, and future without converting them
into three absolute instants. The present is the locus of coupling between two
orientations; past and future are senses of transport. Route inversion
exchanges both senses and leaves the coupling point as a section
of boundary.

\subsection{TPK as a transport system}

TPK brings together support, orientation, and memory. Its objects are
finite or coinductive prepared sections; its morphisms are recorded routes
compatible with the readers. Each route contains
\[
\Gamma=(\text{origin},\text{target},\text{sheet},\text{orientation},
\text{carry},\text{boundary},\text{record}).
\]
Composition concatenates paths and composes memories. Inversion reverses the
order and orientation. The readers, the \(108\)-phase calendar, boundary, and
record are internal operators, not added external observers.

A published operation \(f:X\to Y\) admits a lift
\[
 \widetilde f:X\longrightarrow Y\times M_f,
 \qquad \pi_Y\circ\widetilde f=f.
\]
Given fixed lifts of \(f\) and \(g\), the composition has an
associated canonical lift
\[
 \widetilde{g\circ f}:X\longrightarrow Z\times M_f\times M_g.
\]
Memory is not a comment on the operation: it is part of its enriched codomain.

\begin{proposition}[Genealogy of the operations]
In TPK, sum, product, power, root, logarithm, derivation, and integration
admit a common reading:
\begin{center}
\begin{tabular}{@{}p{.18\textwidth}p{.31\textwidth}p{.38\textwidth}@{}}
\toprule
Operation & Action & Preserved Memory \\
\midrule
Sum & aggregation of alternatives & origin and carry \\
Product & stage composition & order and folded fiber \\
Power & monoid iteration & depth and orbit \\
Root & multivalued inversion & branch and orientation \\
Logarithm & iteration coordinate & generator and caliber \\
Derivation & differ by scale & border and orientation \\
Integration & composition of increments & path and boundary conditions \\
\bottomrule
\end{tabular}
\end{center}
\end{proposition}

Exponentiation does not introduce an operation ontologically disconnected from addition and
product: it iterates composition. The root inverts that iteration and must preserve
the branch selection. The logarithm records depth. This hierarchy explains
why nonadic powers and their vacancies are informative for APP: the raised
operation preserves layers of memory that an isolated digital root
contracts.

\subsection{Temporal periodicity of (108) phases and return laws at (27), (54), and (108) iterations}

The observable calendar has three distinct thresholds. After 27 iterations,
the position returns and the orientation is reversed; after 54, position
and orientation are restored; after 108, the complete observable phase is restored. The memory
of the positive semigroup, however, is not erased. Schematically,
\[
F_{\rm obs}^{27}(x,o)=(x,-o),\qquad
F_{\rm obs}^{54}(x,o)=(x,o),\qquad
F_{\rm obs}^{108}(x,o,\phi)=(x,o,\phi).
\]
This hierarchical order will be crucial for separating outward propagation, return,
orientation, and complete cycle. A finite observable period can coexist with
nonperiodic genealogical memory.

\begin{falsadorBox}[title={Type controls}] It is inadmissible to identify TPK with a list of 34 operators; to call the structure a “fibered category” without Cartesian lifts; to place the carry simultaneously as a cellular coboundary and as a 2-cocycle without separating their complexes; or to convert the block \(SO^+(1,1)\) into physical spacetime before constructing the metric representation.
\end{falsadorBox}


\section{Hopf fibration, Villarceau circles, double covering, and Möbius holonomy}

\subsection{Complex structure and Hopf fibration}

Let \(S^3\subset\mathbb C^2\) and
\[
h(z_0,z_1)=
\bigl(2\Re(z_0\overline z_1),
2\Im(z_0\overline z_1),|z_0|^2-|z_1|^2\bigr)
\in S^2.
\]
On the Hopf torus
\[
T_\eta=\{(z_0,z_1):|z_0|=\cos\eta,
|z_1|=\sin\eta\},
\]
the curves
\[
\gamma_+(t)=(\cos\eta\,e^{it},\sin\eta\,e^{it}),
\]
\[
\gamma_-(t)=(\cos\eta\,e^{it},\sin\eta\,e^{-it})
\]
have distinct behaviors: \(\gamma_+\) is a Hopf fiber and
\(\gamma_-\) projects with degree two. Under stereographic projection, the
homology classes \((1,1)\) and \((1,-1)\) produce the two diagonal
Villarceau families.

The total complex conjugation
\[
K(z_0,z_1)=(\overline z_0,\overline z_1)
\]
preserves each unoriented family and reverses its orientation. It does not interchange
the two families. Partial conjugation can interchange them, but it is neither a
complex-linear map nor the standard global antiunitary. This correction
is essential to avoid manufacturing a particle--antiparticle symmetry from the
toroidal drawing.

\subsection{Double covering and spinor return}

The map \(SU(2)\to SO(3)\) is a double cover. A loop of
\(2\pi\) in \(SO(3)\) can lift to \(-I\) in \(SU(2)\); the loop of
\(4\pi\) recovers \(I\). The sign structure does not turn a circle into a
Möbius band. Obtaining a band requires a real line bundle
whose holonomy character is \(-1\).

Let \(L_\chi\to S^1\) be the associated bundle to
\[
\chi:\pi_1(S^1)\longrightarrow\{\pm1\},
\qquad \chi(1)=-1.
\]
Its total space is a Möbius band. The route \(C_M=-\operatorname{rev}\)
provides the sign candidate; the Villarceau circle provides the loop; the
physical intertwiner must prove that the former is the holonomy of the latter.

\subsection{Spectral decomposition of the oriented memory operator and derivation of the ratio (5{:}2)}

With \(B_c=7I+2S\), define the oriented memory operator
\[
B_{\rm mem}^{\varepsilon}=5I+2\varepsilon S,
\qquad\varepsilon=\pm1.
\]
For \(\varepsilon=+1\),
\[
\Spec(B_{\rm mem}^{+})=\{7,3\},
\qquad\det B_{\rm mem}^{+}=21.
\]
Interpreting the values \((7,3)\) as the outer and inner equatorial boundaries
in a common positive chart yields
\[
R+r=7\ell,\qquad R-r=3\ell,
\]
and therefore
\[
R=5\ell,\qquad r=2\ell,\qquad
\sin\beta=\frac rR=\frac25.
\]
The Villarceau angle is
\[
\beta=\arcsin\frac25
=23.5781784782\ldots^\circ.
\]

The algebraic theorem \(\{7,3\}\leftrightarrow(5,2)\) is exact. The geometric
premise interpreting the spectrum as the two boundaries of a torus is a
declared interface. It must not be concealed, nor may the angle be used as a target.

The realization of the full block \(B_c\) retains another reading:
\(R:r=7:2\) and \(\beta=\arcsin(2/7)\). These are not rival values; they belong to
two functors: total block and memory mode. The physical priority of the latter
must be decided by the intertwiner with the dynamics, not by decimal proximity to
another constant.

\subsection{Conjugate Angular Coordinates and Phase Unit Normalization in the Hopf Fibration}

HMT’s native angle is a phase \(u\in\mathbb R/\mathbb Z\). The readers
publish
\[
\theta_{\rm rad}=2\pi u,
\qquad
\theta_{\rm deg}=360u.
\]
Radians and degrees are distinct charts of the same circular element. The radian is not more “ontological” by virtue of being dimensionless; its naturality follows from arc length. Degrees provide another absolute coordinate of the full turn. Any relation with \(\alpha^{-1}\), \(\varphi\), or CKM must first be formulated as an invariant quotient or holonomy and only thereafter evaluated in an angular chart.

The exact identity already available
\[
\tanh^2(2\log\varphi)=\frac59
\]
connects the Perron--Fibonacci channel to the spectral quotient of \(B_c\). Its
complement is \(4/9\), the normalized gap. This relation is structural.
It does not imply that every angle constructed with \(\varphi\) is a physical angle.

\subsection{Brouwer, Möbius, and degree}

Brouwer's fixed-point theorem, fractional linear Mobius transformations,
and the Mobius band are three distinct objects. They are related only
after fixing domain, action, and bundle. The elliptic advance and
retardation orbits may be modeled as two sections of a bundle with sign
holonomy; the Mobius transformation describes the conformal action in the chart;
Brouwer's theorem guarantees a fixed point for continuous maps of the disk.
None of these statements replaces the realization diagram.

\begin{fronteraBox}[title={Result and boundary}]
The degree-two character of the Hopf curve, the orientation of the
total conjugation, and the algebraic construction \(7,3\mapsto5,2\) are proved. The
toroidal reading (5{:}2), Möbius holonomy, and the
particle--antiparticle interpretation form a chain of interfaces that must share a
single intertwiner. The volume does not force that intertwiner through an
angular coincidence.
\end{fronteraBox}


\section{Finite Hilbertian realization of the path system and its conditional projectors}

\subsection{Matrix realization of each depth}

For \(d\ge1\), define
\[
  \mathcal C_9^{(d)}=(\mathbb Z/9\mathbb Z)^d,
  \qquad
  H_d=\ell^2(\mathcal C_9^{(d)}),
  \qquad
  A_d=B(H_d)\cong M_{9^d}(\mathbb C).
\]
The canonical basis \(\{\delta_x:x\in\mathcal C_9^{(d)}\}\) preserves the
discrete label of each state. The uniform measure induces the normalized
trace
\[
  \tau_d(a)=\frac{1}{9^d}\Tr(a).
\]

The factorization
\[
  \mathcal C_9^{(d+1)}
  =
  \mathcal C_9^{(d)}\times\mathbb Z/9\mathbb Z
\]
produce a unitary identification
\[
  H_{d+1}\cong H_d\otimes\mathbb C^9.
\]
This is the exact input to the operator tower. This does not identify
\((\mathbb Z/9)^3\) with \(\mathbb F_3^6\) as groups: they share \(729\)
elements, but their operations and carry cocycles are distinct.

\subsection{Weyl–Heisenberg Representation of Phase and Displacement on the Triadic Fiber}

Let \(\omega=e^{2\pi i/9}\). On
\(H_1=\ell^2(\mathbb Z/9\mathbb Z)\), we define
\[
  (Xf)(r)=f(r-1),
  \qquad
  (Zf)(r)=\omega^r f(r).
\]
Then
\[
  X^9=Z^9=I,
  \qquad
  ZX=\omega XZ.
\]
The group generated by \(X,Z,\omega I\) is the finite Heisenberg group of
order \(9^3=729\).

\begin{theorem}[Stone--von Neumann finite]
Every irreducible unitary representation of the finite Heisenberg group in
which the center acts by means of the primitive character
\(\omega I\) is unitarily equivalent to the preceding representation.
\end{theorem}

\begin{proof}
Let \(v\) be an eigenvector of \(Z\). The vectors
\[
  v,Xv,\ldots,X^8v
\]
have eigenvalues distinct from \(Z\), since
\(ZX^kv=\omega^kX^kZv\). They are orthogonal and span an invariant
subspace. By irreducibility, they form a basis of the entire representation.
In that basis, \(X\) is the cyclic shift and \(Z\) the diagonal phase.
\end{proof}

\begin{controlbox}
Uniqueness fails if the primitive central character is not fixed or if
reducible representations are allowed. This finite theorem must not be confused
with the continuous version for the canonical commutation relations.
\end{controlbox}

\subsection{Conditional projectors and pre-Heisenberg incompatibility}

The additive and multiplicative sheets of \(\APP\) define information
subspaces and their orthogonal projectors
\(P_{\Sigma,d},P_{\Pi,d}\). The constructed identities satisfy
\[
 \|[P_{\Sigma,2},P_{\Pi,2}]\|
 =\frac{\sqrt2}{3},
 \qquad
 \|[P_{\Sigma,3},P_{\Pi,3}]\|
 =\frac{\sqrt3}{4}.
\]
Noncommutativity is not deduced from a quantum analogy: it is a finite
matrix theorem. Its "pre-Heisenberg" reading is subsequent to that calculation.

\begin{proposition}
The preceding projectors do not jointly belong to the same
Boolean subalgebra of projections.
\end{proposition}

\begin{proof}
In an abelian subalgebra all projectors commute. The displayed commutators
are nonzero.
\end{proof}

\subsection{Hilbertian structure of the matrix publication and preserved genealogical memory}

Hilbert space provides linear superposition, orthogonal projection,
and spectral calculus. HMT supplies the genealogical index of the basis and distinguishes
which operator represents a transition, a measurement, or a shadow. The
realization does not replace \(\TPK\): it represents it.

\begin{sourcebox}
The APP projectors and the Weyl--Heisenberg pair are part of the preceding
construction. The organization \(H_d,A_d\) and the uniqueness proof are established
here by applying the corresponding classical theorem to the HMT
objects.
\end{sourcebox}


\section{Finite symplectic structure of APP and nonadic fibration of observables}
\label{p12-83-c2-s4:gmc:chap:simplectica}

Symplectic geometry is not part of the published theorems on
De las Cuevas--Netzer magic squares. The legitimate link passes through a
neighboring, standard route: Weyl operators, discrete phase space and
mutually unbiased bases \citep{gibbons2004phase,planat2007rings}. In HMT,
that route acquires a nonadic lift and fiber memory.

\subsection{The alternating form and the Weyl commutation}

On $V_3=\mathbb F_3^2$ we define
\[
 \langle(q,p),(q',p')\rangle_{\rm sp}=qp'-pq'.
 \tag{GMC-5.1}
\]
If $X|k\rangle=|k+1\rangle$,
$Z|k\rangle=\omega^k|k\rangle$, and $\omega=e^{2\pi i/3}$, the operators
$W_{q,p}=X^qZ^p$ satisfy, under these conventions,
\[
 W_uW_v=\omega^{-\langle u,v\rangle_{\rm sp}}W_vW_u.
 \tag{GMC-5.2}
\]
Therefore, two operators commute exactly when their labels are
orthogonal for the alternating form. In dimension two, the maximal
isotropic sublines are the four points of
\[
 \mathbb P^1(\mathbb F_3)=\{[1:0],[0:1],[1:1],[1:2]\}.
 \tag{GMC-5.3}
\]
Each generates a maximal abelian context of $M_3(\mathbb C)$; its
spectral projectors constitute a \PVM. Contexts associated with distinct
lines produce the four MUB bases of \textup{(GMC-2.13)}.

\begin{proposition}[Qutrit decomposition by directions]
\label{p12-83-c2-s4:gmc:prop:descomposicion-masa-qutrit}
Let $\mathcal D_a$ be the four diagonal subalgebras determined by the
lines of \textup{(GMC-5.3)} and $\mathcal D_a^0$ their traceless parts. Then
\[
 M_3(\mathbb C)
 =\mathbb CI\oplus
   \bigoplus_{a\in\mathbb P^1(\mathbb F_3)}\mathcal D_a^0
 \tag{GMC-5.4}
\]
as an orthogonal sum for the Hilbert--Schmidt product. The dimension
count is $9=1+4\cdot2$.
\end{proposition}

\begin{proof}
Each $\mathcal D_a^0$ has dimension two. Unbiasedness implies
$\Tr(A^*B)=0$ for $A\in\mathcal D_a^0$ and
$B\in\mathcal D_b^0$, $a\ne b$. All are orthogonal to $I$. The
dimensions sum to nine, the dimension of $M_3(\mathbb C)$.
\end{proof}

This identity explains why the twelve projections are not twelve arbitrary
independent pieces: they are four resolutions of the identity whose centered
parts decompose the entire qutrit operator space.

\subsection{The lift \texorpdfstring{$12\to4\times3$}{12 to 4 by 3}}

On the modulo nonadic $V_9=(\mathbb Z/9\mathbb Z)^2$ it preserves the forma
alternante
\[
 \langle(a,b),(c,d)\rangle_9=ad-bc\pmod9.
 \tag{GMC-5.5}
\]
The projective line $\mathbb P^1(\mathbb Z/9\mathbb Z)$ consists of
the primitive vectors modulo multiplication by a unit. It can
be represented as
\[
 \{[1:t]:t\in\mathbb Z/9\mathbb Z\}
 \ \sqcup\ 
 \{[3u:1]:u\in\mathbb Z/3\mathbb Z\},
 \tag{GMC-5.6}
\]
and has $9+3=12$ points. Reduction modulo three gives the diagram
\[
 \mathbb P^1(\mathbb Z/9\mathbb Z)
 \xrightarrow{\ r\ }
 \mathbb P^1(\mathbb F_3),
 \qquad |r^{-1}(a)|=3.
 \tag{GMC-5.7}
\]

The basis $a$ of a qutrit context and the outcome $k$
of its \PVM likewise form a set $4\times3$. A
\emph{calibration} is a bijection
\[
 \Xi:\mathbb P^1(\mathbb Z/9\mathbb Z)
 \longrightarrow
 \{P_{a,k}:a\in\mathbb P^1(\mathbb F_3),\ k\in\mathbb F_3\}
 \tag{GMC-5.8}
\]
that sends each fiber of $r$ to the three outcomes of the context lying over
$a$.

\begin{controlbox}
The fibration is canonical; assigning the internal order of each fiber to the
three outputs $k$ is a gauge. A bijection is not yet an
equivariance. To promote $\Xi$ to an intertwiner, one must prove that the
relevant linear, projective, and antiunitary actions are transported in
every fiber. The preliminary gauge fails an antiunitary check in one
fiber; that failure is preserved as a falsifier, not concealed.
\end{controlbox}

\subsection{Depth, carry, and memory over the discrete phase space}

The Weyl geometry over $\mathbb F_3^2$ explains commutation, contexts, and
MUB. HMT adds three levels that are not recovered from the symplectic form by
itself:

\begin{enumerate}
 \item the lift of a ternary direction to three points of
 $\mathbb P^1(\mathbb Z/9\mathbb Z)$;
 \item the sheet, carry and orientation that distinguish states with the same reduction;
 \item the TPK route that transports those coordinates and can return to the
 phase without returning to the state.
\end{enumerate}

The potential contribution to the qutrit literature is not “another construction of
four MUBs,” which is already known. It is a geometry of \emph{context roots}: each
ternary context has three nonadic antecedents, and those antecedents
may carry dynamics and monodromy. The order-twelve square of
\Cref{p12-83-c22-s3:gmc:thm:qms-12-3} is the first normalized shadow of that atlas.

\subsection{Typed separation between the qutrit overlap
\texorpdfstring{$1/3$}{1/3} and the HMT reader
\texorpdfstring{$1/\pi$}{1/pi}}

In a pair of MUB qutrit bases,
\[
 \Tr(P_{a,k}P_{b,\ell})=\frac13.
 \tag{GMC-5.9}
\]
In the HMT reader defined here, the relation between a real sheet and a
volumetric reading can take the quotient
\[
 \frac{\pi}{\pi^2}=\frac1\pi.
 \tag{GMC-5.10}
\]
The two expressions are not interchangeable: $1/3$ is an overlap between
projectors in dimension three; $1/\pi$ is a scalar weight of another reader.
The difference does not prevent their coordination. For a fixed \PVM
$\{P_0,P_1,P_2\}$, we define a measurement with nine effects:
\[
 F_k^{(\pi)}=\frac3\pi P_k\quad(k=0,1,2),
 \qquad
 F_r^{(\pi)}=\frac{1-3/\pi}{6}I_3\quad(r=3,\ldots,8).
 \tag{GMC-5.11}
\]
It is a \POVM because $3/\pi<1$ and
\[
 \sum_{r=0}^{8}F_r^{(\pi)}
 =\frac3\pi I_3+\left(1-\frac3\pi\right)I_3=I_3.
 \tag{GMC-5.12}
\]
If $Q$ is a rank-one projector of a distinct MUB context,
\[
 \Tr(QF_k^{(\pi)})
 =\frac3\pi\Tr(QP_k)
 =\frac3\pi\cdot\frac13
 =\frac1\pi.
 \tag{GMC-5.13}
\]

\begin{bridgebox}
The formula \textup{(GMC-5.11)} exactly transforms the MUB probability $1/3$ into the reader $1/\pi$ through the gauge $3/\pi$, without identifying the two numbers. The six residual effects preserve the unit and complete a \POVM of nine outcomes, from which another cyclic \QMS $(9,3)$ is obtained. This construction does not yet identify the geometric provenance of the residual mass; it fixes the exact operator type that it must have.
\end{bridgebox}

\subsection{Symplectic group and equivariance equation}

In dimension two over $\mathbb F_3$, the symplectic group coincides with
$SL(2,\mathbb F_3)$ and permutes the four lines. Its action lifts
projectively, by means of the Clifford group, on the \PVM qutrits. Let
$G_{\rm HMT}$ be a group or groupoid of transformations of the nonadic atlas and
$U_g$ the candidate unitary or antiunitary operator in the qutrit chart. The
condition that turns a calibration into a dynamical bridge is
\[
 \Xi(g\cdot\ell)=U_g\,\Xi(\ell)\,U_g^*,
 \qquad
 g\in G_{\rm HMT},\quad
 \ell\in\mathbb P^1(\mathbb Z/9\mathbb Z).
\tag{GMC-5.14}
\]
If $g$ is an arrow and not merely a global symmetry, the equation must include
the source and target and is formulated as a natural transformation. The verification of
\textup{(GMC-5.14)} is the step that separates a fortunate labeling from a
structural transduction.

\begin{sourcebox}
The symplectic and MUB identities rely on the finite-phase literature
\citep{gibbons2004phase,planat2007rings}. They are not attributed to the papers on
magic squares by De las Cuevas--Netzer. The nonadic fibration, calibration,
and route readers constitute the original contribution of the
HMT construction studied here.
\end{sourcebox}
